\documentclass[11pt]{article}
\usepackage[a4paper,margin=25mm]{geometry}
\usepackage{amsmath,amssymb,amsthm,hyperref}
\newtheorem{theorem}{Theorem}
\title{An identity-jump counterexample to 00000008963}
\author{Ziang Ni (\texttt{ziangni-sys})}
\date{10 October 2026}
\begin{document}
\maketitle
The conjecture describes Riemann--Hilbert reconstruction from jump
matrices and asserts that its Fredholm solvability criterion is the
\emph{vanishing} of a determinant. The Chinese text explicitly says
that the determinant equals zero. With the standard Fredholm
convention $D=\det(I-C)$, the identity-jump problem is a counterexample.

\paragraph{The actual normalized Riemann--Hilbert problem.}
Let $\Sigma=\{z\in\mathbb C:|z|=1\}$, oriented counterclockwise, and
choose the scalar jump $J(z)=1$. This is the $1\times1$ matrix case.
Seek a function $Y$ holomorphic off $\Sigma$, with boundary values
satisfying $Y_+(z)=Y_-(z)J(z)$, and normalized by $Y(z)\to1$ at
infinity. The actual function
\[
 Y(z)=1\quad(z\in\mathbb C)
\]
is entire, has both boundary values equal to $1$, satisfies the jump,
and has the required normalization. Thus the problem is solvable.

\paragraph{The actual correction operator and determinant.}
The usual singular-integral equation uses
\[
 C_{J-I}f=C_-\bigl(f(J-I)\bigr),\qquad (I-C_{J-I})\mu=I.
\]
Here $f(J-I)$ is identically zero for \emph{every} function $f$.
Its Cauchy transform, side limits and principal value are therefore
zero. This is a calculation of the actual operator, not an assumption
of an invertibility certificate. For example, parameterizing the circle
by $s(t)=e^{it}$ gives the Cauchy transform
\[
 \frac1{2\pi i}\int_0^{2\pi}
  \frac{f(s(t))(J(s(t))-1)}{s(t)-z}\,i s(t)\,dt=0
 \qquad (z\notin\Sigma).
\]
Consequently $C_{J-I}=0$, a rank-zero trace-class operator, and
$I-C_{J-I}=I$ on the boundary-function space.

For a finite-rank correction $C$, the standard Fredholm determinant
is the determinant of $I-C$ restricted to any finite-dimensional
invariant subspace containing the range of $C$. The range here is
$\{0\}$, so on any such $n$-dimensional subspace the matrix is
literally $I_n-0$. Thus, independently of the chosen dimension,
\[
 D=\det(I-C_{J-I})=\det(I_n)=1\ne0.
\]
This is also immediate from the Fredholm series: the zeroth term is
$1$ and every positive-order minor of the zero kernel vanishes.
In particular the finite-rank reduction is the standard Fredholm
determinant, rather than a newly chosen scalar invariant.

\begin{theorem}
The stated vanishing criterion in 00000008963 is false.
\end{theorem}
\begin{proof}
The normalized identity-jump problem above has an explicit full
solution, while its standard Fredholm determinant equals $1$.
Hence solvability does not imply determinant zero.
\end{proof}

\paragraph{Convention and Lean verification.}
This counterexample uses the conventional Fredholm determinant
$\det(I-C)$; it does not rule out artificially defining some different
quantity that vanishes on solvable data. The familiar invertibility
criterion uses \emph{nonvanishing}, consistent with this example.
The Lean project verifies complex differentiability, genuine boundary
continuity and jump matching, and the constant limit along the
cobounded filter (normalization at infinity). It computes the actual
parametrized Cauchy integral with the jump factor, proves that the
resulting linear operator is zero, and calculates $\det(I_n-0)=1$
for every finite dimension. Its final theorem combines solvability
with the nonzero determinant. Lean 4.19.0 and mathlib commit
\texttt{c44e0c8ee63ca166450922a373c7409c5d26b00b} are pinned.
All printed axiom audits contain only the permitted standard axioms.
\end{document}
